\documentclass[11pt]{amsart}

\usepackage[margin=1.15in]{geometry}
\usepackage{amsmath,amssymb,amsthm}
\usepackage{graphicx}
\usepackage{booktabs}
\usepackage{xcolor}
\definecolor{linkblue}{rgb}{0.1,0.1,0.45}
\usepackage[colorlinks=true,linkcolor=linkblue,citecolor=linkblue,urlcolor=linkblue]{hyperref}

\theoremstyle{plain}
\newtheorem{theorem}{Theorem}[section]
\newtheorem{proposition}[theorem]{Proposition}
\newtheorem{lemma}[theorem]{Lemma}
\newtheorem{corollary}[theorem]{Corollary}
\newtheorem{conjecture}[theorem]{Conjecture}
\newtheorem{question}[theorem]{Question}
\newtheorem{mainthm}{Theorem}
\renewcommand{\themainthm}{\Alph{mainthm}}

\theoremstyle{definition}
\newtheorem{definition}[theorem]{Definition}
\newtheorem{problem}[theorem]{Problem}

\theoremstyle{remark}
\newtheorem{remark}[theorem]{Remark}
\newtheorem{example}[theorem]{Example}

\newcommand{\HH}{\mathbb{H}}
\newcommand{\CC}{\mathbb{C}}
\newcommand{\RR}{\mathbb{R}}
\newcommand{\ZZ}{\mathbb{Z}}
\newcommand{\NN}{\mathbb{N}}
\newcommand{\CP}{\mathbb{CP}}
\newcommand{\HP}{\mathbb{HP}}
\newcommand{\RP}{\mathbb{RP}}
\renewcommand{\SS}{\mathbb{S}}
\newcommand{\re}{\operatorname{Re}}
\newcommand{\im}{\operatorname{Im}}
\newcommand{\id}{\operatorname{id}}
\newcommand{\Aut}{\operatorname{Aut}}
\newcommand{\eE}{\mathrm{e}}
\newcommand{\EE}{\mathcal{E}}
\newcommand{\dd}{\,d}
\DeclareMathOperator{\osc}{osc}

\begin{document}

\title[The zero-sphere program and the Second Involution Conjecture]{The zero-sphere program for the Riemann Hypothesis\\ and the Second Involution Conjecture}

\author{Edward B. Baker III}
\email{edwardbaker86@gmail.com}

\date{August 2026}

\begin{abstract}
We describe a research program on the Riemann Hypothesis organized around the
canonical extension of the completed zeta function $\Xi$ to the quaternions as an intrinsic slice
regular function. The zero set of the extension consists of real points and round $2$-spheres, and RH
is equivalent to the absence of spheres. We prove the program's structure theorems (a twistor lift to
an incidence quadric on which RH becomes an equivariant statement about the components of a divisor),
its two no-go theorems (no continuous linear functional of a rotational extension, and no
deformation-invariant degree count, detects spheres), and its numerically certified explicit-formula
growth criterion. The program culminates in the \emph{Second Involution Conjecture}: the
zero multiset carries the holomorphic involution $\rho\mapsto1-\rho$ and the antiholomorphic
involution $\rho\mapsto\bar\rho$, RH says they agree, and the conjecture asserts that the
conjugation-paired ``reality defect'' is approximable by a semialgebraic function of finitely many
prime-computable explicit-formula samples at finite logarithmic prime scale. The final sections prove
that this conjecture, exactly as first formulated, is \emph{true unconditionally --- but for a
degenerate reason}: its adaptive sampling quantifier admits an intermediate-value-theorem encoding
that writes the answer into the choice of two sample points, consuming no arithmetic (the mechanism
works verbatim for any Euler-product-free zero multiset whose window observable oscillates). The
proof's one substantive ingredient is of independent interest: an unconditional
uniform-oscillation lemma for the prime-computable observable $c\mapsto\re Z(c,\log2)$, with
effective constants at a natural window length. We diagnose the loophole, prove a
Kronecker-recurrence lower bound showing that the non-adaptive version of the conjecture has genuine
content (continuous semialgebraic estimators at scale $\Lambda$ cannot resolve an isolated defect at
height $x_0$ unless $\Lambda\ge(1-o(1))\log\log x_0$), and restate the conjecture in the
class-uniform form on which the program now rests, with the shifted family
$\Xi_\alpha(z)=\Xi(z+i\alpha)\Xi(z-i\alpha)$ as first testbed. This is a program paper: it claims no
progress on RH itself.
\end{abstract}

\maketitle
\setcounter{tocdepth}{1}
\tableofcontents

%=====================================================================
\section{Introduction}\label{sec:intro}

Rotate the critical line onto the real axis: with
\[
\xi(s)=\tfrac12 s(s-1)\pi^{-s/2}\Gamma(s/2)\zeta(s),\qquad
\Xi(z):=\xi\!\left(\tfrac12+iz\right),
\]
$\Xi$ is entire, even, of order $1$, real on $\RR$, and the Riemann Hypothesis says all its zeros
are real. Because the Taylor coefficients of $\Xi$ are real, $\Xi$ extends canonically to the
quaternions $\HH$ as an \emph{intrinsic slice regular} function, and by the zero-structure theorem
for such functions its zero set is a union of isolated real points and isolated round $2$-spheres:
real zeros of the complex $\Xi$ give points, conjugate pairs $a\pm iy_0$ give spheres $a+y_0\SS$.
\textbf{RH holds if and only if the quaternionic $\Xi$ has no zero spheres} (Theorem~\ref{thm:A}).

This paper is a self-contained account of the program built on that dichotomy, together with a
complete treatment of its newest development. The program's foundational results --- two structure
theorems, two no-go theorems, one certified growth criterion --- are proved in
\S\S\ref{sec:reformulation}--\ref{sec:bridge}, and the numerical laboratory certifying every
computational claim is included as ancillary code (\S\ref{sec:numerics}). The paper's second half
(\S\S\ref{sec:conjecture}--\ref{sec:repair}) states the program's organizing conjecture --- the
\emph{Second Involution Conjecture} --- and then proves three things about it.

\begin{enumerate}
\item \emph{The conjecture as first formulated is true, unconditionally --- by a loophole}
(Theorem~\ref{thm:encoding}). Its sampling quantifier lets the sample points be chosen after the
target height, and an intermediate-value-theorem selection then encodes the answer into the choice
of two points. The mechanism consumes no arithmetic and applies verbatim to any Euler-product-free
zero multiset whose window observable oscillates; it passes the conjecture's own hardness checks
vacuously. The result is a diagnosis, not
progress: a specification with an existential quantifier over continuum-valued sample parameters,
ranging after the instance, is satisfiable by a sealed envelope.
\item \emph{The proof's one substantive ingredient survives as unconditional infrastructure}
(Lemma~\ref{lem:osc}): the prime-computable explicit-formula observable $F(c)=\re Z(c,\log2)$ has
uniform positive oscillation on every window of every fixed length, at every height, with
effective constants at a natural window length --- the prime-$2$ term is never drowned by the
archimedean drift.
\item \emph{The repaired, non-adaptive conjecture has genuine content}
(Theorem~\ref{thm:lower}): in a sampling model abstracting exactly the structure of finite-scale
prime data, non-adaptive continuous semialgebraic estimators at scale $\Lambda$ cannot resolve an
isolated defect at height $x_0$ unless $\Lambda\ge(1-o(1))\log\log x_0$. Adaptivity had been
silently projecting this wall to zero.
\end{enumerate}

The paper ends with the two repaired formulations --- the class-uniform
Conjecture~\ref{conj:uniform}, which we adopt as the program's headline, and the single-function
non-adaptive form Conjecture~\ref{conj:repaired} --- and with a sharply posed, RH-free first test: compute
the smoothed paired zero count from finite-scale prime data for the shifted family
$\Xi_\alpha(z)=\Xi(z+i\alpha)\Xi(z-i\alpha)$, whose reality defect is nonzero, unbounded, and
exactly decomposable.

\subsection*{What this paper is not}
It is not a proof strategy with a claimed endgame, and Theorem~\ref{thm:encoding} is not a proof of
a standing open conjecture in any useful sense: it shows the original statement was formalized at
the wrong joint, salvages its unconditional content, and repairs it. The no-go theorems and the two
hardness checks (Davenport--Heilbronn; de Bruijn--Newman) run through everything precisely to
prevent over-reading, and \S\ref{sec:bridge} states plainly where the program currently reduces to
classical difficulty. The reader is entitled to keep score against the standing question
\emph{``what does the reformulation buy?''}

\subsection*{Provenance}
The program synthesizes a sequence of single-author manuscripts (2015--2026) --- a quaternionic
extension hierarchy for $\xi'/\xi$ with an $SU(2)$-representation-theoretic decomposition and a
negative result on linear pairings \cite{Baker2015,Baker2016}, and twistor-geometric machinery
developed for generalized conformal maps and Yang--Mills symmetries
\cite{Baker2017,Baker2022arXiv,Baker2023arXiv} --- with the modern theory of slice regular
functions (Gentili--Struppa \cite{GentiliStruppa2007} and its development
\cite{GentiliStoppato2008,CGSS2009,GSSbook,GhiloniPerotti2011}), whose intrinsic subclass is
exactly the function class the manuscripts were searching for, and whose twistor branch
\cite{GSS2014,Altavilla2018,AltavillaSarfatti2019} realizes the intended geometry. The historical
manuscripts are archived unaltered at \cite{ZenodoArchive}. The synthesis, the computations, the
numerical laboratory, the formulation of the Second Involution Conjecture, and the resolution and
repair reported in \S\S\ref{sec:resolution}--\ref{sec:repair} were carried out in August 2026 with
substantial AI research assistance; see the Acknowledgements for a precise statement.

%=====================================================================
\section{The reformulation}\label{sec:reformulation}

Let $\SS=\{J\in\HH:J^2=-1\}$ be the $2$-sphere of imaginary units; each $J\in\SS$ spans a slice
$\CC_J=\RR+J\RR\cong\CC$. A function on an axially symmetric slice domain is \emph{slice regular}
\cite{GentiliStruppa2007} if it is holomorphic on every slice, and \emph{intrinsic} if its power
series has real coefficients; equivalently $f(x+yJ)=\alpha(x,y)+J\beta(x,y)$ with
$F:=\alpha+i\beta$ holomorphic and real on $\RR$ (the \emph{stem function}; the extension
$F\mapsto f$ is the rotation map of \cite[eq.~(2.5)]{Baker2015}). Intrinsic functions form a
commutative real algebra under \emph{pointwise} operations, restricting on each slice to the
ordinary complex algebra, and are determined by one slice via the representation formula
\cite{CGSS2009,GSSbook}. The zero set of a nonzero slice regular function is a union of isolated
points and isolated spheres $x+y\SS$ (Gentili--Stoppato \cite{GentiliStoppato2008}).

Riemann's integral representation extends verbatim: with $\Phi$ the positive, even, rapidly
decaying kernel of \cite[\S10.1]{Titchmarsh1986},
\begin{equation}\label{eq:cosine}
\Xi(q)=4\int_0^\infty\Phi(u)\cos(qu)\dd u=\sum_{n\ge0}c_{2n}q^{2n},\qquad q\in\HH,
\end{equation}
absolutely convergent on all of $\HH$ since $\cos(qu)$ is entire intrinsic and $\Phi$ decays like
$\eE^{-\pi\eE^{2u}}$; the positivity and decay are due to P\'olya \cite{Polya1926}. The
normalizing constant $4$ depends on conventions for $\Phi$; with the convention displayed it is
certified numerically (\S\ref{sec:numerics}: the identity holds to ten digits against an
independent computation of $\xi(\tfrac12+ix)$, and the right side vanishes to the quadrature
floor at the first zero ordinate).

\begin{mainthm}[Reformulation]\label{thm:A}
RH holds if and only if the zero set of $\Xi:\HH\to\HH$ defined by \eqref{eq:cosine} contains no
$2$-sphere.
\end{mainthm}

\begin{proof}
Immediate from the zero-structure theorem and the stem-function dictionary: nonreal zeros of the
complex $\Xi$ occur in conjugate pairs $a\pm iy_0$ (reality on $\RR$) and correspond exactly to the
zero spheres $a+y_0\SS$ of the extension.
\end{proof}

By itself this is a relabeling; its value is the structures it makes available, and everything
below is an attempt to convert the relabeling into leverage.

%=====================================================================
\section{Twistor geometry: the equivariant dichotomy}\label{sec:twistor}

Inside the twistor space $\CP^3\to\HP^1\supset\HH$ of $S^4$, with real structure $\sigma$ (the
fixed-point-free antiholomorphic involution whose invariant lines are the twistor fibers), sits the
smooth \emph{incidence quadric}
\[
Q=\{(q,\eta):\eta\in\SS,\ q\in\CC_\eta\}\ \cong\ \CP^1_\zeta\times\CP^1_w ,
\]
swept by the lines $L_\zeta=\{(x+\eta y,\eta)\}$, $\zeta=x+iy$ --- the tautological lifts of the
spheres $x+y\SS$, degenerating to twistor fibers over $\RR$. Here $w\in\CC\cup\{\infty\}$
parametrizes the sphere of complex structures stereographically,
$\eta(w)=(w+j)^{-1}\,i\,(w+j)\in\SS$, so that the antipodal map $\eta\mapsto-\eta$ is
$w\mapsto-1/\bar w$; the point of $Q$ with coordinates $(\zeta,w)$ is $q=x+y\,\eta(w)$. The two
rulings are the sphere-lifts ($\zeta$ fixed) and the fixed-complex-structure slices ($w$ fixed);
$\sigma$ restricts to $\rho(\zeta,w)=(\bar\zeta,-1/\bar w)$, conjugation times the antipodal
map, again fixed-point free.

\begin{mainthm}[Lift]\label{thm:B}
Let $f$ be intrinsic slice regular on an axially symmetric slice domain $\Omega$, with stem
function $F$, and let $\Omega_Q\subseteq Q$ be the part of the quadric over $\Omega$. Then the
product map $\tilde f=F\times\id:\Omega_Q\to Q$ is holomorphic and $\rho$-equivariant, and
covers $f$: over the point $q=x+y\,\eta(w)$ it returns
$f(q)=\alpha(x,y)+\beta(x,y)\,\eta(w)$ with the same complex-structure label $w$. If $f$ is
entire, $\tilde f$ is defined on all of $Q$ minus the ruling at infinity.
\end{mainthm}

\begin{proof}
Holomorphy is immediate: $F$ is holomorphic in $\zeta$ by intrinsicness and the second factor is
the identity. Equivariance is the computation
$\tilde f(\rho(\zeta,w))=(F(\bar\zeta),-1/\bar w)=(\overline{F(\zeta)},-1/\bar w)
=\rho(\tilde f(\zeta,w))$, using reality of $F$ on $\RR$. Covering: the point $(\zeta,w)\in Q$ is
$q=x+y\,\eta(w)$ where $\zeta=x+iy$; writing $F(\zeta)=\alpha(x,y)+i\beta(x,y)$, the image point
$(F(\zeta),w)$ is $\alpha+\beta\,\eta(w)$, which equals $f(q)$ by the intrinsic form
$f(x+yJ)=\alpha+J\beta$ applied in the slice $\CC_{\eta(w)}$.
\end{proof}

\begin{remark}[Relation to the Gentili--Salamon--Stoppato lift]\label{rem:GSS}
For intrinsic $f$, Theorem~\ref{thm:B} is the degenerate product form of the
Gentili--Salamon--Stoppato twistor lift \cite[Thm.~5.3 and eq.~(5.6)]{GSS2014} (substitute
$h=0$, $\hat g=g$ there); the ruled geometry of such lifts is developed in
\cite{Altavilla2018,AltavillaSarfatti2019}. What we add is the systematic use of this product
form: the identification of its natural domain and target as the incidence quadric $(Q,\rho)$
with its two rulings, the equivariance, the maximality statements below, and the divisor-level
dichotomy of Theorem~\ref{thm:C}. We have not found these statements in the literature, but the
ingredients are all close to \cite{GSS2014,Altavilla2018,AltavillaSarfatti2019}, and we make no
strong novelty claim.
\end{remark}

\subsection*{Maximality of the quadric}
In what sense is $Q$ the natural home of the lift? Care is needed: $Q$ is a smooth hypersurface,
so unconstrained local holomorphic extensions off $Q$ always exist and the interesting question
fixes the twistor structure. Two statements result.

\begin{proposition}[The na\"ive extension is holomorphic only along $Q$]\label{prop:naive}
Let $f$ be intrinsic slice regular and not affine ($f\not\equiv aq+b$, $a,b\in\RR$), and consider
the fiber-preserving assignment $\Psi_f:(q,\eta)\mapsto(f(q),\eta)$ on the full correspondence
space $\HH\times\SS$, with the complex structure at $(q,\eta)$ given by left multiplication by
$\eta$ on the $\HH$-factor and the standard one on the $\SS$-factor. Then the Cauchy--Riemann
system for $\Psi_f$ holds at every point of $Q$ (where it reduces to slice holomorphy), and on no
open subset of $(\HH\times\SS)\setminus Q$.
\end{proposition}

\begin{proof}
Holomorphy at $(q,\eta)$ in the $\HH$-directions requires $d_qf(\eta v)=\eta\,d_qf(v)$ for all
$v\in\HH$. Decompose $T_q\HH=\CC_J\oplus\CC_J^\perp$ where $q=x+yJ$, $y>0$. For a monomial
$f=q^m$: on the slice plane $d_qf(v)=mq^{m-1}v$ (slice holomorphy), while for $v\perp\CC_J$,
using $vq=\bar qv$,
\[
d_qf(v)=\Bigl(\sum_{i=0}^{m-1}q^i\bar q^{\,m-1-i}\Bigr)v
=\frac{q^m-\bar q^m}{q-\bar q}\,v
\]
(the sum collapses to the divided difference because $q\bar q=|q|^2$ is real). By
$\RR$-linearity and locally uniform convergence, for any slice regular $f$,
\begin{equation}\label{eq:blocks}
d_qf(v)=f'(q)\,v \quad (v\in\CC_J),\qquad
d_qf(v)=\partial_s f(q)\,v \quad (v\in\CC_J^\perp),
\end{equation}
where $f'$ is the slice derivative and
$\partial_sf(q)=(f(q)-f(\bar q))(q-\bar q)^{-1}$ is the \emph{spherical derivative}
\cite{GhiloniPerotti2011,GSSbook}; both act by left multiplication by elements of $\CC_J$. Take
$\eta\notin\{\pm J\}$ and split $\eta=\eta_\parallel+\eta_\perp$, $\eta_\perp\neq0$. Testing the
intertwining condition on $v\in\CC_J$, the slice parts agree ($\CC_J$ is commutative) and the
orthogonal parts require $\partial_sf(q)\,\eta_\perp v=\eta_\perp f'(q)v$, i.e.\ (since
$\eta_\perp w=\bar w\eta_\perp$ for $w\in\CC_J$) the pointwise identity
$\partial_sf(q)=\overline{f'(q)}$. For intrinsic $f$ with stem $F=\alpha+i\beta$ the spherical
derivative $\beta(x,y)/y$ is \emph{real}, so on any open set where the identity holds, $F'$ is
real-valued, hence constant (open mapping theorem), hence $F$ is affine with real coefficients
--- excluded. So the identity fails on a dense open subset of every neighborhood off $Q$. On $Q$
itself ($\eta=\pm J$) the intertwining condition is exactly slice holomorphy, and the fiber
directions impose nothing.
\end{proof}

\begin{proposition}[Rigidity of quadric immersions]\label{prop:immersion}
Every holomorphic immersion $\Phi:\CP^1\times\CP^1\to\CP^3$ is the Segre embedding composed with
projective and M\"obius transformations (possibly after the factor swap); in particular the image
is a smooth quadric and $\Aut(Q)$ consists of restrictions of projective transformations.
\end{proposition}

\begin{proof}
Write $S=\CP^1\times\CP^1$, $\Phi^*\mathcal O(1)=\mathcal O(a,b)$. The image spans $\CP^3$ (an
immersion of a compact surface into a hyperplane $\CP^2$ would be a covering, forcing
$S\cong\CP^2$), and immersivity forces $a,b\ge1$. The complete system $|\mathcal O(a,b)|$ embeds
$S$ into $\CP^N$, $N=(a+1)(b+1)-1$, and $\Phi$ is this embedding followed by a projection from a
center $\Lambda\cong\CP^{N-4}$; the projection is an immersion iff $\Lambda$ avoids the tangent
variety $\mathrm{Tan}(S)$, which has dimension $4$ (surfaces with degenerate tangent varieties
are classified; see \cite{GriffithsHarris1979} or \cite[Ch.~I]{Zak1993}). Since
$\dim\Lambda+\dim\mathrm{Tan}(S)=N$, the center always meets $\mathrm{Tan}(S)$ unless $N=3$,
i.e.\ $(a,b)=(1,1)$: $\Phi$ is the Segre embedding up to $PGL_4\times(PGL_2\times PGL_2)$. For
the last claim, any automorphism of $Q$ preserves $\mathcal O(1,1)$ (the square root of $-K_Q$ in
the torsion-free Picard group), by which $Q$ is embedded, so it extends to $\CP^3$. Classical
background: \cite[Ch.~6]{GriffithsHarris1978}.
\end{proof}

\begin{remark}\label{rem:maximality}
Proposition~\ref{prop:naive} says the incidence relation $q\in\CC_\eta$ is forced: the lift
cannot be prolonged off $Q$ compatibly with the fibration by complex structures unless $f$ is
M\"obius. Proposition~\ref{prop:immersion} complements it globally: there is no ``bigger''
holomorphic model of the intrinsic class inside twistor space (by Riemann--Hurwitz,
$\tilde f$ composed with the Segre embedding is an immersion iff $F$ is M\"obius). Together they
resolve the rigidity tension of \cite{Baker2017} --- fiberwise degree-one self-maps of all of
twistor space are M\"obius --- by relocating the rotational function class from twistor space to
the incidence quadric. A version of Proposition~\ref{prop:immersion} may well be classical; we
could not locate a reference in this form, and the proof is assembled from classical facts.
\end{remark}

\subsection*{The equivariant dichotomy}
The zero divisor of $\tilde f$ in $Q$ is $Z(\tilde f)=F^{-1}(0)\times\CP^1_w$, a union of
first-ruling lines with multiplicities, one per zero of $F$.

\begin{mainthm}[Equivariant dichotomy]\label{thm:C}
Let $f\not\equiv0$ be intrinsic slice regular, with zero divisor
$Z(\tilde f)=\sum_km_kL_{\zeta_k}$. Then:
\textup{(1)} a real zero of $f$ contributes a $\sigma$-\emph{fixed} component (a twistor fiber)
and a zero sphere a $\sigma$-\emph{swapped pair} of components, the multiset $\{\zeta_k\}$ being
$\sigma$-invariant with multiplicity;
\textup{(2)} in $H_2(Q;\ZZ)$ a real zero of multiplicity $m$ contributes $m[\mathrm{ruling}]$
and a sphere of multiplicity $m$ contributes $2m[\mathrm{ruling}]$ --- a pair of simple real
zeros and a simple sphere have the same class, so homology cannot separate the two cases;
\textup{(3)} the quotient $Y=Q/\sigma$ is a smooth compact $4$-manifold on which real zeros
descend to $\RP^2$ components of the zero locus and zero spheres to $S^2$ components.
Consequently, for $f=\Xi$:
\emph{RH holds iff every component of the zero locus of $\tilde\Xi$ in $Y$ is nonorientable.}
\end{mainthm}

\begin{proof}
(1) $\sigma(L_\zeta)=L_{\bar\zeta}$ directly from $\rho(\zeta,w)=(\bar\zeta,-1/\bar w)$ (for
real $\zeta$ the line is $\sigma$-invariant with the antipodal map $w\mapsto-1/\bar w$ acting on
the parameter); the multiset is conjugation-closed with multiplicity because $F$ is real on
$\RR$. (2) In $H_2(Q;\ZZ)=\ZZ[\mathrm{A}]\oplus\ZZ[\mathrm{B}]$ each $L_\zeta$ has class
$[\mathrm{A}]$; a real zero of multiplicity $m$ gives $mL_{x_0}$, a sphere the pair
$mL_{a+iy_0}+mL_{a-iy_0}$. (3) $\sigma$ is fixed-point free on $\CP^3$, hence on $Q$, so $Y$ is
smooth and $Q\to Y$ is a double cover; $\rho$ covers $\zeta\mapsto\bar\zeta$, so $Y$ fibers over
the closed half-plane $\{\im\zeta\ge0\}\cup\{\infty\}$, with interior fiber the $2$-sphere
$\CP^1_w$ (the swapped pair $L_\zeta,L_{\bar\zeta}$ identified into one) and boundary fiber
$L_x/(\text{antipodal})\cong\RP^2$. The displayed equivalence is Theorem~\ref{thm:A} read in
$Y$.
\end{proof}

%=====================================================================
\section{Two no-go theorems}\label{sec:nogo}

\begin{mainthm}[Linear no-go]\label{thm:D}
Let $E$ be any continuous $\RR$-linear ``rotational extension'' map from the intrinsic entire
functions (locally uniform topology) to a topological vector space of $\HH$-valued functions ---
e.g.\ the slice extension, or any Fueter-type extension.
Then every continuous linear functional of $E(F)$ is a continuous linear functional of $F$ itself,
and no continuous linear functional detects reality of zeros: no linear pairing of any rotational
extension --- against any test function, over any contour, in any sector --- separates on-axis from
off-axis zeros.
\end{mainthm}

\begin{proof}
The first claim is the composition of continuous linear maps. For the second, suppose a
continuous linear $\Lambda_0$ and a threshold $c$ cut out the reality locus on the two-parameter
family $F_{t,s}(\zeta)=(\zeta^2+t)(\zeta^2+s)$, whose zeros are all real iff $t\le0$ and $s\le0$
(a quarter-plane $P$). Expanding, $F_{t,s}=\zeta^4+(t+s)\zeta^2+ts$, so by linearity
$\Lambda_0(F_{t,s})=A+(t+s)B+ts\,C$ with $A=\Lambda_0(\zeta^4)$, $B=\Lambda_0(\zeta^2)$,
$C=\Lambda_0(1)$. If the reality locus were $\{\Lambda_0\ge c\}$: along the boundary ray
$\{t=0,\,s\le0\}$, points with $t=0^-$ are inside $P$ and $t=0^+$ outside, so continuity forces
$A+sB=c$ for all $s\le0$, hence $B=0$, $A=c$; the condition then reads $ts\,C\ge0$, whose
solution set is $\{ts\ge0\}$, $\{ts\le0\}$, or everything --- never the quarter-plane. If it
were $\{\Lambda_0=c\}$: the level set of an affine function of $(t+s,ts)$ is empty, a line, all
of $\RR^2$, a hyperbola, or a pair of crossed lines --- never the boundary of a quarter-plane.
This formalizes and explains the negative result of \cite[\S4]{Baker2015}, where the concrete
instance ($E=$ Fueter extension, $\Lambda=$ Cauchy--Fueter pairing) was verified by explicit
Laurent-expansion computation.
\end{proof}

\begin{mainthm}[Degree no-go]\label{thm:E}
For intrinsic slice regular $f$, the real differential at $q=x+yJ$ ($y>0$) splits as left
multiplication by the slice derivative $f'(q)$ on $\CC_J$ and by the spherical derivative
$\partial_sf(q)=(f(q)-f(\bar q))(q-\bar q)^{-1}$ on $\CC_J^\perp$, so
$\det(d_qf)=|f'(q)|^2|\partial_sf(q)|^2\ge0$. The local Brouwer degree of $f/|f|$ (the
$4$-dimensional argument principle) is $m$ at a real zero of multiplicity $m$ and $2m$ at a zero
sphere of multiplicity $m$ --- exactly the complex count: deformation-invariant integer counts
cannot detect spheres.
\end{mainthm}

\begin{proof}
The block decomposition is \eqref{eq:blocks}, established in the proof of
Proposition~\ref{prop:naive}; left multiplication by $a\in\HH$ on a $2$-plane it preserves acts
with determinant $|a|^2$, giving the product formula for $\det(d_qf)$.

\emph{Real zero.} For $f=q^m$ on $|q|=1$, a generic target value has exactly the $m$ slice roots
as preimages, each positively oriented by the determinant formula (on the orthogonal plane the
divided difference is $\sin m\psi/\sin\psi\cdot(\text{unit})$ for $q=e^{J\psi}$, positive at
generic points); hence local degree $m$. A general intrinsic $f$ with a real zero of multiplicity
$m$ factors locally as $(q-x_0)^m\tilde f(q)$ with $\tilde f(x_0)\neq0$ real, and the degree is
stable under this perturbation.

\emph{Zero sphere.} For the model $f=q^2+1$ near the sphere $\SS$: the family $q^2+t$,
$t\in[0,1]$, is nonvanishing on a fixed large sphere $|q|=R>1$, so the degree there is constant
in $t$; at $t=0$ it equals the degree of $q^2$, which is $2$ by the real-zero computation (a real
zero of multiplicity $2$); and at $t=1$ the zero set is $\SS$. In
$H_3(\HH\setminus\SS)\cong\ZZ$ the large sphere is homologous to the boundary $S^2\times S^1$ of
a tubular neighborhood of $\SS$ (the region between them is a $4$-chain avoiding $\SS$), and the
map is orientation-preserving where nondegenerate, so the local degree at the zero sphere is
$2$; multiplicity $m$ gives $2m$ by the same argument applied to
$(q^2+1)^m\cdot(\text{unit})$.
\end{proof}

\begin{remark}
The count ``real zeros plus twice spheres'' is exactly what the complex argument principle
produces for the even $\Xi$. This is required by the hardness checks of \S\ref{sec:hardness}:
the degree is a deformation invariant and spheres are destroyed by arbitrarily small
deformations, so the degree \emph{must not} see them. The equivariant component type of
Theorem~\ref{thm:C} is the minimal refinement that does --- and it is precisely not
deformation-invariant: components change type exactly when zeros collide on the axis.
\end{remark}

Any successful discriminant must therefore be nonlinear \emph{and} finer than homology and degree
--- equivariant (\S\ref{sec:twistor}), positivity-type (\S\ref{sec:positivity}), or spectral.

%=====================================================================
\section{Hardness checks}\label{sec:hardness}

Two tests are run against every candidate mechanism in the program.

\emph{Davenport--Heilbronn.} Let $\chi$ be the odd Dirichlet character mod $5$ with $\chi(2)=i$,
and set
\[
\kappa=\frac{\sqrt{10-2\sqrt5}-2}{\sqrt5-1},\quad
f_{DH}(s)=\tfrac12\bigl[(1-i\kappa)L(s,\chi)+(1+i\kappa)L(s,\bar\chi)\bigr],\quad
\Lambda_{DH}(s)=\Bigl(\tfrac{5}{\pi}\Bigr)^{s/2}\Gamma\Bigl(\tfrac{s+1}{2}\Bigr)f_{DH}(s).
\]
The function $\Xi_{DH}(z)=\Lambda_{DH}(\tfrac12+iz)$
\cite{DavenportHeilbronn1936a,DavenportHeilbronn1936b,Titchmarsh1986} is entire of order $1$, even,
real on $\RR$ --- so it has an intrinsic slice regular extension with exactly the symmetry of $\Xi$
--- but has zeros off the critical line and no Euler product. Its extension has genuine zero
spheres; the first three, recomputed here by Newton iteration at $25$-digit precision
(\S\ref{sec:numerics}; matching Spira \cite{Spira1994}), have centers and radii
\begin{equation}\label{eq:DHspheres}
(a,y_0)=(85.699348,\,0.308517),\quad(114.163343,\,0.150830),\quad(166.479306,\,0.074356).
\end{equation}
\textbf{Requirement 1:} symmetry plus growth cannot exclude spheres; a candidate mechanism must
consume arithmetic input and must demonstrably fail for $\Xi_{DH}$.

\emph{De Bruijn--Newman.} By Rodgers--Tao \cite{RodgersTao2020} the de Bruijn--Newman constant is
nonnegative: under backward heat flow the zeros leave the axis immediately, so $\Xi$ sits at the
exact boundary of the sphere-free regime, with spheres of arbitrarily small radius appearing under
arbitrarily small deformation. \textbf{Requirement 2:} any true criterion degrades only
polynomially (or worse) as the sphere radius $y_0\to0$; a soft argument stable under small
perturbations proves too much and is wrong.

%=====================================================================
\section{The positivity hierarchy}\label{sec:positivity}

On the real axis, $-(\log\Xi)''(x)=\sum_\rho(x-\rho)^{-2}$: real zeros contribute positively, and a
sphere $(a,y_0)$ contributes negatively \emph{exactly on} $|x-a|<y_0$, its equatorial shadow. The
Laguerre inequality $L_1[\Xi]=\Xi'^2-\Xi\Xi''\ge0$ is thus the dichotomy's bottom certificate, and
the full generalized Laguerre family characterizes Laguerre--P\'olya membership --- hence RH ---
by Csordas--Varga \cite[Thm.~2.9]{CsordasVarga1990} (necessity) and Csordas--Vishnyakova
\cite[Thm.~2.3]{CsordasVishnyakova2013} (sufficiency, with no order or genus restriction).
Numerically (\S\ref{sec:numerics}): $L_1[\Xi]>0$ with relative margin $0.173$ on $[0,52]$, while
$L_1[\Xi_{DH}]<0$ precisely inside the predicted windows $|x-a|<y_0$ at the computed spheres.

The grading that separates sphere radii is by spherical Bessel tags: the \emph{sector energies}
\begin{equation}\label{eq:sector}
\EE_\ell=\int_0^\infty\Bigl|\sum_\rho j_\ell(\lambda y_\rho)\,\eE^{i\lambda a_\rho}\Bigr|^2
w(\lambda)\dd\lambda,\qquad \rho=a_\rho+iy_\rho,\ y_\rho\ge0
\end{equation}
(each conjugate pair counted once, by its upper member), are quadratic in the zeros, so
Theorem~\ref{thm:D} does not apply, and they are graded:

\begin{lemma}[Grading]\label{lem:grading}
Let the zero multiset be finite (in practice: windowed) and closed under conjugation, and let
$w>0$ be continuous with all moments finite. Then \textup{(1)} real zeros contribute nothing to
any $\ell\ge1$ sector ($j_\ell(0)=0$); \textup{(2)} a single sphere of radius $y_0$ contributes
to $\EE_\ell$ at order $y_0^{2\ell}$ as $y_0\to0$, with leading coefficient
$((2\ell+1)!!)^{-2}\int\lambda^{2\ell}w\dd\lambda$; \textup{(3)} all zeros are real $\iff$
$\EE_\ell=0$ for all $\ell\ge1$ $\iff$ $\EE_1=0$.
\end{lemma}

\begin{proof}
(1) and (2) are immediate from $j_\ell(t)=t^\ell/(2\ell+1)!!\,(1+O(t^2))$
\cite[10.52.1]{DLMF}. For (3): if $\EE_1=0$ then, since $w>0$, the finite sum
$S(\lambda)=\sum_\rho j_1(\lambda y_\rho)\eE^{i\lambda a_\rho}$ vanishes identically; as
$\lambda\to0^+$, $S(\lambda)=\tfrac\lambda3\sum_\rho y_\rho+O(\lambda^2)$, so
$\sum_\rho y_\rho=0$, and since every $y_\rho\ge0$, all vanish.
\end{proof}

The tags are not ad hoc: they arise as the radial profiles of the $SU(2)$-isotypic decomposition
of the rotational extension developed in the program's underlying manuscripts
\cite{Baker2015,Baker2016} (archived at \cite{ZenodoArchive}), where the image of the exponential
in the $\ell$-th sector is $i^\ell(j_\ell(\lambda y)+i\,j_{\ell+1}(\lambda y))\eE^{\lambda x}$:
$j_\ell(\lambda y)$ is the profile with which a zero at height $y$ enters sector $\ell$ at
frequency $\lambda$. We stress that this grounding is motivation only --- \eqref{eq:sector} is
used as a \emph{definition} here, and nothing below depends on the unpublished material.
Numerically (\S\ref{sec:numerics}): measured log--log slopes of $\EE_\ell$ versus $y_0$ are
$2.000$, $4.000$, $6.000$ for $\ell=1,2,3$; the DH spheres fit the $y_0^2$ law in $\ell=1$ to
$2\%$.

%=====================================================================
\section{The explicit-formula bridge}\label{sec:bridge}

Rotated coordinates throughout: $\rho=\tfrac12+i\gamma$, the zero multiset $\{\gamma\}$ closed
under $\gamma\mapsto-\gamma$ and $\gamma\mapsto\bar\gamma$, $|\im\gamma|<\tfrac12$. For $h$ even,
holomorphic and decaying on $|\im r|\le\tfrac12+\varepsilon$, with
$g(u)=\tfrac1{2\pi}\int h(r)\eE^{-iur}\dd r$, the Guinand--Weil explicit formula
\cite{Guinand1948,Weil1952}, in the normalization of \cite[Thm.~5.12]{IwaniecKowalski2004} rotated,
reads
\begin{equation}\label{eq:EF}
\sum_\gamma h(\gamma)=h\bigl(\tfrac i2\bigr)+h\bigl(-\tfrac i2\bigr)
+\frac1{2\pi}\int_{-\infty}^\infty h(t)\Bigl[\re\,\psi\bigl(\tfrac14+\tfrac{it}2\bigr)
-\log\pi\Bigr]\dd t
-2\sum_{n\ge2}\frac{\Lambda(n)}{\sqrt n}\,g(\log n),
\end{equation}
certified numerically as stated (\S\ref{sec:numerics}): on the Gaussian family below, the zero
side computed from $250$ zeros matches pole${}+{}$archimedean${}+{}$prime side to
$10^{-8}$--$10^{-10}$ at five $(c,\lambda)$ configurations, and the $\sigma=1$,
$\lambda^*=\log2$ implementation of \S\ref{sec:resolution} agrees to $2.7\times10^{-14}$. Fix a width $\sigma>0$, set
$G(u)=\eE^{-u^2/2\sigma^2}$, and take the test family and linear building block
\begin{equation}\label{eq:Z}
h_{c,\lambda}(r)=\eE^{i\lambda(r-c)}G(r-c)+\eE^{-i\lambda(r+c)}G(r+c),
\qquad
Z(c,\lambda)=\sum_\gamma h_{c,\lambda}(\gamma).
\end{equation}
By \eqref{eq:EF}, $Z(c,\lambda)$ is \emph{unconditionally computable} from
$\{\Lambda(n):|\log n-\lambda|=O(1/\sigma)\}$ plus explicit archimedean terms: the frequency
$\lambda$ is the logarithmic prime scale.

\begin{lemma}[Diagonal anomaly]\label{lem:anomaly}
Fix $c>0$ and $\sigma$, and write $Y=\sup_\rho y_\rho$ for a zero multiset with
$y_\rho\le\tfrac12$ and $N(T)\ll T\log T$. \textup{(1)} A real zero pair $\{\pm a\}$ contributes
a bounded oscillatory amount to $Z(c,\lambda)$, uniformly in $\lambda$. \textup{(2)} A zero
sphere $(a,y_0)$ --- the quadruple $\{\pm a\pm iy_0\}$ --- contributes to $|Z(c,\lambda)|^2$ a
diagonal block oscillating, for large $\lambda$, between $4(G\mp G')^2\eE^{2\lambda y_0}$ up to
$O(\eE^{\lambda y_0})$, where $G=\eE^{-((a-c)^2-y_0^2)/2\sigma^2}$,
$G'=\eE^{-((a+c)^2-y_0^2)/2\sigma^2}$ (note $G>G'$ for $a,c>0$): exponential growth at rate exactly
$2y_0$ in the $\limsup$ sense, weighted by the Gaussian distance of the sphere's center from the
window. \textup{(3)} Unconditionally $|Z(c,\lambda)|\le C_{c,\sigma}\,\eE^{\lambda Y}$, and the
total of all products in $|Z|^2$ with at least one factor from a zero of height
$\le Y-\varepsilon$ is $O(\eE^{\lambda(2Y-\varepsilon)})$.
\end{lemma}

\begin{proof}
(1) For real $\gamma$, $|h_{c,\lambda}(\gamma)|\le G(\gamma-c)+G(\gamma+c)\le2$. (2) Since
$h_{c,\lambda}$ is even, the quadruple contributes $2h(\gamma)+2h(\bar\gamma)$,
$\gamma=a+iy_0$; the growing terms are the first window term of $h(\bar\gamma)$,
$\eE^{i\lambda(a-c)}\eE^{\lambda y_0}G\eE^{i\theta}$ (with $\theta=(a-c)y_0/\sigma^2$ from the
complex Gaussian), and the second window term of $h(\gamma)$,
$\eE^{-i\lambda(a+c)}\eE^{\lambda y_0}G'\eE^{i\theta'}$; all other terms have modulus
$\eE^{-\lambda y_0}O(1)$, and the squared modulus of the sum oscillates between the stated
bounds. (3) Each $\gamma_k=a_k+iy_k$ satisfies
$|h(\gamma_k)|\le2\eE^{\lambda y_k}\eE^{y_k^2/2\sigma^2}\eE^{-(a_k-c)^2/2\sigma^2}$; with
$y_k\le\tfrac12$ the middle factor is bounded, and summing the Gaussian against the zero-counting
measure of density $O(\log(2+|a_k|))$ gives the bound; the refined estimate uses
$\eE^{\lambda y_k}\le\eE^{\lambda(Y-\varepsilon)}$ on the low factor of each such product.
\end{proof}

\begin{mainthm}[Rate criterion]\label{thm:F}
Let $Y=\sup_\rho y_\rho$ for a zero multiset satisfying the hypotheses of
Lemma~\ref{lem:anomaly}.
\textup{(a)} Unconditionally,
$\limsup_{\lambda\to\infty}\lambda^{-1}\log|Z(c,\lambda)|\le Y$ for every window; in particular RH
implies $Z$ bounded in $\lambda$ in every window, so exponential growth in any single window
disproves RH. \textup{(b)} If either $Y=0$, or $Y>0$ and for some $\varepsilon>0$ only finitely
many zeros have $y_\rho>Y-\varepsilon$ (a non-accumulation hypothesis; the $Y=0$ case is listed
separately since under RH every zero satisfies $y_\rho>-\varepsilon$ and the finiteness clause
is vacuously false), then for every fixed window
\[
\limsup_{\lambda\to\infty}\frac{\log|Z(c,\lambda)|}{\lambda}=Y=\Theta-\tfrac12,
\]
with $\Theta$ the supremum of real parts of the zeros; under this hypothesis RH is equivalent to
subexponential growth of $Z$ in every window, and a positive rate reads off the supremal sphere
radius exactly.
\end{mainthm}

\begin{proof}
(a) is Lemma~\ref{lem:anomaly}(3). (b) If $Y=0$, part (a) gives $\limsup\le0$, and the
mean-value argument below, applied to the bounded real-zero sum with $\lambda$-independent
amplitudes, bounds $|Z|$ below along a subsequence, so the $\limsup$ is exactly $0$. If $Y>0$:
fix $\varepsilon$ as given, let $\gamma_1,\dots,\gamma_K$ (with mirrors) be the finitely many
zeros of height $>Y-\varepsilon$, and split $Z=Z_{\mathrm{top}}+Z_{\mathrm{rest}}$ with
$|Z_{\mathrm{rest}}|\le C\eE^{\lambda(Y-\varepsilon)}$ by Lemma~\ref{lem:anomaly}(3). By part
(2) of the lemma, $\eE^{-\lambda Y}Z_{\mathrm{top}}$ is, up to $o(1)$, a finite exponential sum
$\sum_kd_k\eE^{-\lambda(Y-y_k)}\eE^{i\lambda b_k}$ with distinct frequencies after grouping
(coincident frequencies group coherently: the growing contributions at frequency $a_k-c$ carry
the same Gaussian phase). With $y_{\max}=\max_ky_k\in(Y-\varepsilon,Y]$, the Besicovitch mean
value of $|\eE^{-\lambda y_{\max}}Z_{\mathrm{top}}|^2$ over $[\Lambda,2\Lambda]$ tends to
$\sum_{k:y_k=y_{\max}}|d_k|^2>0$, so along a sequence $\lambda_n\to\infty$,
$|Z(c,\lambda_n)|\ge\tfrac{c'}2\eE^{\lambda_n(Y-\varepsilon)}$ eventually; hence
$\limsup\ge Y-\varepsilon$, and shrinking $\varepsilon$ gives equality with (a). The
identification $Y=\Theta-\tfrac12$ is the rotation $s=\tfrac12+iz$.
\end{proof}

\begin{remark}[The non-accumulation gap]\label{rem:accumulation}
The hypothesis in (b) is automatic for Davenport--Heilbronn in any fixed height range and holds
for $\zeta$ under RH (the $Y=0$ branch); under $\neg$RH it would follow if the supremum $\Theta$
is attained, which is not known. Removing it --- heights accumulating at the supremum, centers
escaping to infinity --- requires a multi-scale refinement of the mean-value argument that we
have not carried out; we flag this as a genuine, if technical, gap between (a) and the clean
unconditional equivalence. Direction (a), the falsifiable one, is unconditional.
\end{remark}

On real data (\S\ref{sec:numerics}): $250$ Riemann ordinates give fitted rate $-0.0003$; DH with
its first sphere gives rate $0.3085$ against the true $y_0=0.30852$ --- the radius of a genuine
off-line zero recovered to four decimals from the growth of a single window.

\begin{remark}[Honest accounting]\label{rem:honest}
At the exponential-rate level Theorem~\ref{thm:F} is a Gaussian-mollified form of the classical
equivalence between zero-free regions and prime-counting error terms: detecting a sphere of radius
$y_0$ through growth requires frequencies $\lambda\gtrsim1/y_0$, i.e.\ primes up to
$\eE^{1/y_0}$. The wall is the old wall, with its geometry exposed. Genuinely new leverage, if
any, lies below the exponential scale --- which is what the Second Involution Conjecture is about.
\end{remark}

%=====================================================================
\section{Two involutions, and the Second Involution Conjecture}\label{sec:conjecture}

The zero multiset carries two commuting involutions:
\[
\iota_{\mathrm{FE}}:\rho\mapsto1-\rho\ \ (\gamma\mapsto-\gamma),\qquad
\iota_{\mathrm{R}}:\rho\mapsto\bar\rho\ \ (\gamma\mapsto\bar\gamma),
\]
the first holomorphic (functional equation), the second antiholomorphic (reality; the real
structure $\sigma$ of \S\ref{sec:twistor}). RH is precisely the statement that they coincide on
the zeros. The explicit formula \eqref{eq:EF} computes sums of \emph{analytic} functions of
$\gamma$; since the multiset is conjugation-closed, $\sum_\gamma h(\bar\gamma)=\sum_\gamma
h(\gamma)$ for every such $h$: \emph{the explicit formula natively sees $\iota_{\mathrm{FE}}$ and
is structurally blind to $\iota_{\mathrm{R}}$}. The sharpest scalar illustration: the
functional-equation-paired sum $\sum_\rho\rho^{-1}(1-\rho)^{-1}=2+\gamma_{\mathrm{E}}-\log4\pi$ is
classical \cite[Ch.~12]{Davenport2000}, while the conjugation-paired twin
$\sum_\rho|\rho|^{-2}$ is computed by no known method --- and the comparison is one-sided:

\begin{proposition}\label{prop:gateway}
Unconditionally $\sum_\rho|\rho|^{-2}\ \ge\ \sum_\rho\rho^{-1}(1-\rho)^{-1}
=2+\gamma_{\mathrm{E}}-\log4\pi$, with equality if and only if RH holds.
\end{proposition}

\begin{proof}
Both sums converge absolutely, so group zeros into orbits $\{\rho,\bar\rho,1-\rho,1-\bar\rho\}$ of
the group generated by the two involutions and set $a=1/\rho$, $b=1/(1-\rho)$: a full orbit
contributes $2|a|^2+2|b|^2$ to the left side and $4\re(ab)$ to the right, and
$4\re(ab)\le4|a||b|\le2|a|^2+2|b|^2$, with $|a|=|b|$ iff $|\rho|=|1-\rho|$ iff
$\re\rho=\tfrac12$; on the line the orbit degenerates to $\{\rho,\bar\rho\}$ and the two
contributions agree.
\end{proof}

RH is thus the statement that a single convergent sum achieves its known classical lower bound.
The object the program asks to compute is the \emph{windowed reality defect}
\begin{equation}\label{eq:defect}
D(c)=\sum_\gamma(\im\gamma)^2\,\eE^{-(\re\gamma-c)^2/2\sigma^2}\ \ge\ 0,
\end{equation}
which vanishes identically iff RH holds --- a regularized diagonal of the quadratic class, the
part pairing each zero with its own conjugate. The off-diagonal is smoothed pair correlation
\cite{Montgomery1973,GoldstonMontgomery1987}; the diagonal has no dictionary. In the exact
algebraic model (real polynomials, with coefficients playing the role of primes),
conjugation-paired sums of roots \emph{are} computable from the coefficients --- but only
semialgebraically, never analytically: already for $x^2+bx+c$,
$\sum r_i^2=b^2-2c$ while $\sum|r_i|^2=\max(b^2-2c,\,2c)$, non-analytic across the discriminant.
(Real roots: $\sum|r_i|^2=b^2-2c\ge2c$ since $b^2\ge4c$; a conjugate pair:
$\sum|r_i|^2=2r\bar r=2c$ with $b^2<4c$.) This dictates the shape of the conjecture. With
$\mathcal A(\Lambda)=\{Z(c,\lambda):c\in\RR,\ 0\le\lambda\le\Lambda\}$ the \emph{prime data at
scale $\Lambda$} (every element unconditionally prime-computable by \S\ref{sec:bridge}), the
program's working notes (August 2026; archived at \cite{ZenodoArchive}) first formulated:

\begin{conjecture}[Second Involution Conjecture, original form]\label{conj:original}
For every $\varepsilon>0$ there exist a finite scale $\Lambda(\varepsilon)<\infty$, an integer $N$,
and a semialgebraic function $S$ of $N$ variables, such that for every $c\in\RR$ there are points
$Z(c_1,\lambda_1),\dots,Z(c_N,\lambda_N)\in\mathcal A(\Lambda(\varepsilon))$ with
$|c_j-c|=O(\sigma)$ and
\[
\bigl|\,D(c)-S\bigl(Z(c_1,\lambda_1),\dots,Z(c_N,\lambda_N)\bigr)\bigr|<\varepsilon.
\]
More generally, the same holds for the natural family of conjugation-paired sums
$\sum_\gamma\Psi(\gamma,\bar\gamma)$ with $\Psi$ a divided-difference kernel of exponentials
across the pair $(\gamma,\bar\gamma)$, of which $D(c)$ is the leading example.
\end{conjecture}

The intent was that the defect --- the explicit formula's structural blind spot --- should
nevertheless be a finite-scale, positivity-type functional of prime data, with the quantitative
question (how fast must $\Lambda(\varepsilon)$ grow?) deliberately severed as a separate question.
The next section shows that the statement, as written, failed to capture that intent.

%=====================================================================
\section{Resolution: the conjecture is true as stated, by an encoding}\label{sec:resolution}

\begin{theorem}[Adaptive samples trivialize Conjecture~\ref{conj:original}]\label{thm:encoding}
Fix $\sigma>0$ and set $\lambda^*=\log2$. There exist a constant $r=r(\sigma)>0$ and a
semialgebraic function $S:\CC^2\to\RR$ --- explicitly, with $u:=\re z_1-\re z_2$,
\[
S(z_1,z_2)=
\begin{cases}
\dfrac{u-r/4}{\,r/8-(u-r/4)\,}, & r/4\le u<3r/8,\\[4pt]
0,&\text{otherwise},
\end{cases}
\]
--- such that for every $c\in\RR$ there are points $c_1,c_2\in[c-\sigma,c+\sigma]$, $c_1\ne c_2$,
with
\[
D(c)\;=\;S\bigl(Z(c_1,\log2),\,Z(c_2,\log2)\bigr)\qquad\text{exactly}.
\]
In particular Conjecture~\ref{conj:original} holds unconditionally, with $N=2$, with
$\Lambda(\varepsilon)=\log2$ for every $\varepsilon$, and with error $0$.
\end{theorem}

The only substantive ingredient is an unconditional oscillation lemma for the observable
$F(c):=\re Z(c,\lambda^*)$, which is of independent interest and survives the diagnosis below as
program infrastructure.

\begin{lemma}[Uniform oscillation]\label{lem:osc}
Fix $\sigma>0$, $\lambda^*=\log2$. Then $F(c)=\re Z(c,\lambda^*)$ is real-analytic and even on
$\RR$, and for every window half-length $a>0$ there is $r=r(a,\sigma)>0$ with
\[
\osc_{[c-a,\,c+a]}F\;\ge\;r\qquad\text{for every }c\in\RR .
\]
Moreover the bound is effective at a natural window length: with $\omega=\log2$ and the constants
$B_2,C_0,C_1$ of the decomposition \eqref{eq:decomp} below, for
$a_0=(2/B_2)\bigl(B_2/2\omega+C_1+2C_0/\omega\bigr)$ there is an effective $C_3(\sigma)$ such that
$\osc_{[c-a_0,c+a_0]}F\ge B_2/8$ for all $|c|\ge C_3$ (numerically at $\sigma=1$:
$B_2=0.5407$, $C_0=2.6975$, $C_1=3.7164$, $a_0=43.98$, $B_2/8=0.0676$).
\end{lemma}

\begin{proof}
\emph{Analyticity and decomposition.} Each summand of \eqref{eq:Z} is entire in $c$; uniform
convergence on compacta (Gaussian decay against $N(T)\ll T\log T$) makes $Z(\cdot,\lambda)$ entire,
and closure of $\{\gamma\}$ under conjugation gives $Z(-c,\lambda)=\overline{Z(c,\lambda)}$, so
$F$ is real-analytic and even. Applying \eqref{eq:EF} to $h_{c,\lambda^*}$ and computing the
Gaussian Fourier transform yields the exact decomposition
\begin{equation}\label{eq:decomp}
F(c)=A(c)+E(c)+P(c),\qquad
P(c)=-\sum_{n\ge2}B_n\cos(c\log n),
\end{equation}
where
$B_n=\frac{2\Lambda(n)}{\sqrt n}\cdot\frac{\sigma}{\sqrt{2\pi}}
\bigl(\eE^{-\sigma^2(\log n-\lambda^*)^2/2}+\eE^{-\sigma^2(\log n+\lambda^*)^2/2}\bigr)\ge0$
with $C_0:=\sum_nB_n<\infty$ and $B_2>0$ (this is the Euler product's one appearance), $A$ is the
archimedean term, with $|A'(c)|\le C_A/(1+|c|)$ by the digamma asymptotics, and $E$ is the pole
term, $|E(c)|\le C_E\,\eE^{-c^2/2\sigma^2}$.

\emph{Large $|c|$, effective (the frequency-$\log2$ test).} By evenness take $c\ge C_3$ and
$W=[c-a_0,c+a_0]$. For any $t_0\in W$,
\[
\sup_{t\in W}|P(t)-P(t_0)|\ \ge\ \frac1{2a_0}\Bigl|\int_W\bigl(P(t)-P(t_0)\bigr)\eE^{-i\omega t}\dd
t\Bigr|.
\]
The constant term contributes at most $|P(t_0)|\cdot\frac2\omega\le\frac{2C_0}\omega$. In
$\int_W\cos(t\log n)\eE^{-i\omega t}\dd t$, the $n=2$ resonant term equals $a_0$ plus an error
$\le\frac1{2\omega}$, and the $n\ge3$ terms total at most
$C_1:=\sum_{n\ge3}B_n(|\log(n/2)|^{-1}+(\log2n)^{-1})<\infty$. Hence
$\osc_WP\ge\frac1{2a_0}[B_2a_0-\frac{B_2}{2\omega}-C_1-\frac{2C_0}\omega]=\frac{B_2}4$ by the
choice of $a_0$, and $\osc_WF\ge B_2/4-2a_0C_A/(1+c-a_0)-2C_E\eE^{-(c-a_0)^2/2\sigma^2}\ge B_2/8$
once $c\ge C_3$, $C_3$ effective.

\emph{Arbitrary window length via the Bohr hull.} $P$ has an absolutely convergent almost periodic
expansion, converging uniformly on horizontal strips (since $\sum B_nn^T<\infty$), so $P$ is Bohr
almost periodic and its hull $H(P)$ --- the uniform closure of its translates --- is compact
(Bochner's criterion), consisting of restrictions of entire functions bounded on strips. Every
$Q\in H(P)$ has Fourier--Bohr coefficient of modulus $B_2/2\ne0$ at frequency $\omega$ (the
coefficient functional is continuous under uniform limits), so no $Q$ is constant on any interval:
interval-constancy plus entirety would force constancy on $\RR$ and kill the coefficient. The
oscillation functional $Q\mapsto\osc_{[-a,a]}Q$ is continuous and strictly positive on the compact
$H(P)$, so $r_H(a):=\min_{H(P)}\osc_{[-a,a]}>0$, and
$\osc_{[c-a,c+a]}P\ge r_H(a)$ for every $c$; subtracting the $A,E$ oscillations as above gives
$\osc_{[c-a,c+a]}F\ge r_H(a)/2$ for $|c|\ge C_4(a)$. On the remaining compact range, $F$ is
real-analytic and nonconstant (its oscillation at large $|c|$ is $\ge r_H(a)/2>0$), hence by the
identity theorem nonconstant on every interval, so $\phi(c)=\osc_{[c-a,c+a]}F$ is continuous and
strictly positive, with a positive minimum $r_0$; take $r=\min(r_0,r_H(a)/2)$.
\end{proof}

\begin{remark}\label{rem:deserts}
The constant $r$ collapses (while remaining positive) on windows in zero deserts: numerically at
$\sigma=1$, $\osc_{[-6,6]}F\approx8\times10^{-15}$, while $\osc_{[c-15,c+15]}F\ge2.01$ for all
sampled $c\in[0,1160]\cup\{10^4,10^5,10^6\}$ (\S\ref{sec:numerics}). Lemma~\ref{lem:osc}
quantifies something the program needs anyway: \emph{the prime-$2$ term of the explicit formula is
never drowned by the archimedean drift}, at any height, in any window.
\end{remark}

\begin{proof}[Proof of Theorem~\ref{thm:encoding}]
Apply Lemma~\ref{lem:osc} with $a=\sigma$ to get $r=r(\sigma,\sigma)>0$, and let
$W_c=[c-\sigma,c+\sigma]$, so $\osc_{W_c}F\ge r$ for every $c$. The defect \eqref{eq:defect}
converges, $D(c)\in[0,\infty)$. Define the encoder
$f(d)=\frac r4+\frac r8\cdot\frac d{1+d}$, a semialgebraic increasing bijection of $[0,\infty)$
onto $[r/4,3r/8)$ with inverse $f^{-1}(u)=\frac{u-r/4}{r/8-(u-r/4)}$. Given $c$, put
$u:=f(D(c))$, choose $t_\pm\in W_c$ with $F(t_-)=\min_{W_c}F$, $F(t_+)=\max_{W_c}F$; then
$F(t_+)-F(t_-)\ge r>u>0$, and by the intermediate value theorem there is $c_1$ between $t_-$ and
$t_+$ with $F(c_1)-F(t_-)=u$. Set $c_2:=t_-$. Then
$\re Z(c_1,\log2)-\re Z(c_2,\log2)=u\in[r/4,3r/8)$, so
$S(Z(c_1,\log2),Z(c_2,\log2))=f^{-1}(u)=D(c)$ exactly. The graph of $S$ is a finite Boolean
combination of polynomial equalities and inequalities in the real coordinates, hence
semialgebraic; both samples lie in $\mathcal A(\log2)$; and $\log2$ never depended on
$\varepsilon$. The same argument handles any absolutely convergent real-valued conjugation-paired
sum $\sum_\gamma\Psi(\gamma,\bar\gamma)$ (encode through a semialgebraic homeomorphism of $\RR$
onto a subinterval; complex values coordinatewise), covering the conjecture's ``more generally''
clause.
\end{proof}

\subsection*{Diagnosis}
Theorem~\ref{thm:encoding} must not be mistaken for progress --- not toward RH (nothing here bears
on RH: if RH holds the conjecture is verified by $S\equiv0$ in one line, and the theorem merely
closes the remaining case), and not toward the intended content of the conjecture. It is a
diagnosis, and the two hardness checks of \S\ref{sec:hardness} are the instruments. \emph{First},
the mechanism consumes no arithmetic: the Euler product enters only through $B_2\ne0$ ---
oscillation of one observable --- and the identical encoding ``computes'' the defect of any
two-involutions zero multiset whose window observable oscillates. This is demonstrated numerically
(\S\ref{sec:numerics}) on a synthetic, Euler-product-free multiset carrying an off-line quadruple:
its nonzero defect $0.125$ is recovered to $1.8\times10^{-15}$ from two samples. A formula that
computes the reality defect of Euler-product-free zero sets is not the object the conjecture set
out to specify. \emph{Second}, the conjectured modulus $\Lambda(\varepsilon)$ --- whose growth rate was
the entire quantitative question, with the classical wall at
$\Lambda(\varepsilon)\asymp\varepsilon^{-1/2}$ --- comes out constant, with exact recovery. The
loophole is precisely identifiable: the ``there are points'' quantifier lets the sample points be
chosen after $c$, hence with full knowledge of $D(c)$, and on a window where a continuous
observable has oscillation $\ge r$, a choice of an ordered pair of points transmits an arbitrary
real number in $[0,r]$ --- an information channel of unbounded capacity, before $S$ is consulted.
All information flows through the selection map; none flows through arithmetic. The repair is
not to shrink the window or restrict $S$, but to make the sampling rule non-adaptive.

%=====================================================================
\section{Repair: the non-adaptive and class-uniform conjectures, and a lower bound}\label{sec:repair}

\subsection{The repaired statements}

\begin{conjecture}[Second Involution Conjecture, non-adaptive form]\label{conj:repaired}
Fix $\sigma>0$. For every $\varepsilon>0$ there exist a finite scale $\Lambda(\varepsilon)<\infty$,
an integer $N$, \emph{fixed} offsets $\xi_1,\dots,\xi_N\in[-1,1]$ and \emph{fixed} frequencies
$\lambda_1,\dots,\lambda_N\in[0,\Lambda(\varepsilon)]$, and a semialgebraic function $S$ of $N$
variables, such that
\[
\sup_{c\in\RR}\ \bigl|\,D(c)-S\bigl(Z(c+\xi_1\sigma,\lambda_1),\dots,
Z(c+\xi_N\sigma,\lambda_N)\bigr)\bigr|\;<\;\varepsilon .
\]
\end{conjecture}

The estimator is now one fixed local formula slid along the critical line --- the natural analogue
of a convolution-type explicit formula, restoring the quantifier ordering of Weil positivity (test
objects before zeros) --- and the encoding of Theorem~\ref{thm:encoding} is dead: the samples are
values of $N$ fixed functions of $c$, and any two heights whose data agree must receive the same
estimate. One residual triviality remains: if RH holds, $S\equiv0$ verifies
Conjecture~\ref{conj:repaired}, so the statement is RH-trivial \emph{for $\zeta$ alone}. The
content one actually wants is uniformity over a class in which nonzero defects certifiably occur.
We adopt the following as the program's headline conjecture.

\begin{conjecture}[Second Involution Conjecture, class-uniform form]\label{conj:uniform}
Let $\mathcal X$ be the class of \emph{explicit-formula systems}: pairs $(\Gamma,\beta)$ of a zero
multiset $\Gamma\subset\{|\im\gamma|\le\tfrac12\}$ closed under both involutions with
$N_\Gamma(T)\ll T\log T$, and coefficients $|\beta_n|\le C\log n$, satisfying the Guinand--Weil
identity
\[
\sum_{\gamma\in\Gamma}h_{c,\lambda}(\gamma)=\mathrm{Arch}(c,\lambda)
-2\sum_{n\ge2}\frac{\beta_n}{\sqrt n}\,g_c(\log n)
\]
for all test functions \eqref{eq:Z}, with the fixed archimedean shape of \eqref{eq:decomp} (poles
optional). For every $\varepsilon>0$ there exist $\Lambda(\varepsilon)$, $N$, offsets, frequencies,
and a single semialgebraic $S$ as in Conjecture~\ref{conj:repaired} such that for \emph{every}
$(\Gamma,\beta)\in\mathcal X$,
\[
\sup_{c\in\RR}\ \frac{\bigl|\,D_\Gamma(c)-S\bigl(Z_\Gamma(c+\xi_1\sigma,\lambda_1),\dots\bigr)
\bigr|}{1+D_\Gamma(c)}\;<\;\varepsilon .
\]
\end{conjecture}

Three calibrating observations. \emph{(a)} Davenport--Heilbronn does \emph{not} belong to
$\mathcal X$ as axiomatized: its logarithmic derivative has no absolutely convergent coefficient
side precisely because it has zeros where that series would need to converge. Whether $\mathcal X$
contains members with off-line zeros is therefore not obvious --- and that question (does the
explicit-formula constraint alone force reality?) is a positivity question of exactly
Weil--Bombieri--Connes type \cite{Weil1952,Bombieri2000,Connes2016}; if the constraint does force
reality --- if $\mathcal X$ has no off-line members --- the conjecture collapses to the trivial
$S=0$ and the two-involutions problem is solved outright, which nobody expects to be cheap. \emph{(b)} $\mathcal X$ does contain members with known, nonzero, unbounded
defect: for $0<\alpha<\tfrac12$ the shifted product
$\Xi_\alpha(z)=\Xi(z+i\alpha)\Xi(z-i\alpha)$ lies in a poles-adjusted version of $\mathcal X$ with
$\beta_n=\Lambda(n)(n^\alpha+n^{-\alpha})$, zero multiset $\{\gamma\pm i\alpha\}$, and the exact
identity
\begin{equation}\label{eq:shifted}
D_{\Xi_\alpha}(c)\;=\;2D_\zeta(c)\;+\;2\alpha^2\sum_\gamma\eE^{-(\re\gamma-c)^2/2\sigma^2}.
\end{equation}
\emph{(c)} The relative-accuracy normalization $1+D_\Gamma(c)$ is forced by the lower bound below:
additive accuracy uniform in height is impossible for continuous semialgebraic estimators at any
fixed scale, for reasons independent of RH.

\subsection{The lower bound: non-adaptivity restores content}

We work in a sampling model abstracting exactly the structure established for prime data in the
proof of Lemma~\ref{lem:osc}.

\begin{definition}\label{def:model}
A \emph{data map at scale $\Lambda$} is $V=(V_1,\dots,V_N):\RR\to\CC^N$ with $V_j=A_j+P_j+E_j$,
where each $P_j$ is Bohr almost periodic with absolutely convergent expansion
$P_j(c)=\sum_{n\ge2}b_{j,n}\eE^{\pm ic\log n}$, $\sum_n|b_{j,n}|\le C_P$, and Gaussian truncation
tail $|b_{j,n}|\le C\eE^{-\sigma^2(\log n-\Lambda)^2/2}$ for $\log n\ge\Lambda$; each $A_j$
satisfies $|A_j(c+\tau)-A_j(c)|\le C_A(1+\tau)/(1+|c|)$ and $|A_j(c)|\le C_A\log(2+|c|)$; and
$|E_j(c)|\le C_E\eE^{-|c|/C_E}$. A \emph{defect profile with an isolated quadruple at height
$x_0$} is $D:\RR\to[0,\infty)$ with $D(x_0)\ge d_0>0$ and $D\le d_0/4$ on $[x_0+2\sigma,2x_0]$.
\end{definition}

Non-adaptive samples $V_j(c)=Z(c+\xi_j\sigma,\lambda_j)$, $\lambda_j\le\Lambda$, of any member of
$\mathcal X$ form a data map at scale $\Lambda$; a member whose only off-line zeros are one
conjugate quadruple at $\pm x_0\pm iy_0$ has a defect profile with $d_0=2y_0^2$.

\begin{lemma}[H\"older--growth regularity]\label{lem:holder}
Let $S:\RR^m\to\RR$ be continuous and semialgebraic. There are $C,q>0$ and $\theta\in(0,1]$ such
that $|S(u)-S(v)|\le C(1+R)^q|u-v|^\theta$ for all $|u|,|v|\le R$, every $R\ge1$.
\end{lemma}

\begin{proof}
The modulus $\Phi(R,s)=\sup\{|S(u)-S(v)|:|u|,|v|\le R,\ |u-v|\le s\}$ is finite, continuous and
nondecreasing in $s$ with $\Phi(R,0)=0$, and semialgebraic in $(R,s)$ (first-order definability
plus quantifier elimination). For fixed $R$, the \L ojasiewicz inequality for continuous
semialgebraic functions \cite[Prop.~2.6.4]{BCR} gives $\Phi(R,s)\le C_Rs^{\theta_R}$ on $[0,1]$;
Hardt semialgebraic triviality \cite[Thm.~9.3.2]{BCR} applied to the family $R\mapsto\Phi(R,\cdot)$
shows only finitely many exponents occur, so $\theta=\min_R\theta_R>0$ is uniform; and the optimal
constant $C(R)$ is semialgebraic in $R$, hence polynomially bounded \cite[Prop.~2.6.2]{BCR}. For
$|u-v|\ge1$ use polynomial growth of $|S|$ on balls directly.
\end{proof}

\begin{theorem}[Height-uniformity is impossible at fixed scale]\label{thm:lower}
Fix a data map $V$ at scale $\Lambda$ and a defect profile $D$ with an isolated quadruple at
height $x_0$. If $S:\CC^N\to\RR$ is continuous, semialgebraic, and satisfies
$\sup_{c\in\RR}|D(c)-S(V(c))|\le d_0/8$, then
\[
\Lambda\;\ge\;(1-o(1))\,\log\log x_0 ,
\]
where $o(1)\to0$ as $x_0\to\infty$ for fixed $(d_0,\sigma,N)$ and fixed regularity data
$(C,q,\theta)$ of $S$.
\end{theorem}

\begin{proof}
Assume $\Lambda\le\log\log x_0$. Fix $\delta\in(0,1)$. Truncating each $P_j$ at level $M$ costs
$\sum_{n>M}|b_{j,n}|\le C\int_{\log M}^\infty\eE^{\ell-\sigma^2(\ell-\Lambda)^2/2}\dd\ell
<\delta/4$ upon choosing
$\log M=\Lambda+\sigma^{-2}+\sigma^{-1}\sqrt{2(\Lambda+\log(C''NC/\delta))}$ (complete the square;
the factor $\eE^\ell$ is the integer count per unit of $\log n$). The truncated parts are
trigonometric polynomials on a common frequency set of size $K\le M$. By pigeonhole on the
$K$-torus (partition into $(C'/\delta)^K$ cells; consider the times $\tau_m=2\sigma m$,
$0\le m\le(C'/\delta)^K$), there is $\tau\in[2\sigma,T(\delta)]$, $T(\delta)=2\sigma(C'/\delta)^K$,
with $\|P_j(\cdot+\tau)-P_j\|_\infty<3\delta/4$ for all $j$ (two tail evaluations included).
Provided
\begin{equation}\label{eq:heightcond}
T(\delta)\le\frac{\delta}{4C_A}\,x_0 ,
\end{equation}
we get $\tau\le x_0$, so $x_0+\tau\in[x_0+2\sigma,2x_0]$ and $D(x_0+\tau)\le d_0/4$, while
$|V(x_0+\tau)-V(x_0)|\le2N\delta$ (slow variation of $A_j$, exponential smallness of $E_j$). Both
data values lie in a ball of radius $R=C''\log x_0$, so by Lemma~\ref{lem:holder}
\[
|S(V(x_0))-S(V(x_0+\tau))|\le C(1+C''\log x_0)^q(2N\delta)^\theta\ \le\ d_0/4
\]
for the choice $\delta(x_0)=(\log x_0)^{-2q/\theta}$ and $x_0$ large. Then
$\tfrac{3d_0}4\le|D(x_0)-D(x_0+\tau)|\le\tfrac{d_0}8+\tfrac{d_0}8+\tfrac{d_0}4$, a contradiction.
Hence \eqref{eq:heightcond} fails: $T(\delta(x_0))\ge c\,\delta(x_0)x_0$. Taking logarithms, with
$K\le M=\exp(\Lambda+O_\sigma(\sqrt{\log\log x_0}))$ and $\log(1/\delta)\asymp\log\log x_0$,
\[
\log x_0+O(\log\log x_0)\ \le\ K\log(C'/\delta)\ \le\
\exp\bigl(\Lambda+O_\sigma(\sqrt{\log\log x_0})\bigr)\cdot O(\log\log x_0),
\]
whence $\Lambda\ge\log\log x_0-O_\sigma(\sqrt{\log\log x_0})$.
\end{proof}

\begin{remark}[Caveats and reading]\label{rem:caveats}
(i) Theorem~\ref{thm:lower} constrains \emph{continuous} semialgebraic estimators; a semialgebraic
$S$ may be discontinuous along a lower-dimensional wall, and ruling out wall-crossing requires
controlling where the data curve meets the wall set --- open, though an estimator that changes its
answer discontinuously under $o(1)$ input perturbations violates the de Bruijn--Newman robustness
profile in spirit. (ii) The Dirichlet-type bound $T(\delta)\le(C/\delta)^K$ is wasteful; any
improvement strengthens the conclusion. The point is qualitative: some wall exists for
non-adaptive estimators at every scale, growing with height --- whereas
Theorem~\ref{thm:encoding} shows no wall whatsoever exists for adaptive ones. (iii) Whether the
$\zeta$ multiset realizes the isolated-quadruple hypothesis is exactly as unknowable as RH; the
theorem is a statement about the estimator class over the model --- about what a proof of
Conjectures~\ref{conj:repaired} or \ref{conj:uniform} is allowed to look like. It lives on an axis
orthogonal to the classical exchange-rate wall $\Lambda(\varepsilon)\asymp\varepsilon^{-1/2}$
(defect \emph{size} at given height versus uniformity in \emph{height} at given size); the two
should be read together as the beginning of a two-parameter lower-bound theory that adaptivity had
been projecting to zero.
\end{remark}

\subsection{The first test}
The sharply posed, RH-free entry point is the shifted family: by \eqref{eq:shifted}, proving
Conjecture~\ref{conj:repaired} \emph{for $\{\Xi_\alpha\}$ relative to $\zeta$} is exactly the
problem of computing the smoothed paired count
$\sum_\gamma\eE^{-(\re\gamma-c)^2/2\sigma^2}$ from finite-scale prime data to $o(1)$ --- a
conjugation-paired sum with known nonzero defect, strictly easier than, and on the road to, the
full conjecture. Its cross-term calculus (unconditional and RH-conditional bounds on the
off-diagonal, pair-correlation part of $|Z|^2$ \cite{Montgomery1973,GoldstonMontgomery1987}) is
the program's declared first test, now to be developed with the recurrence adversary of
Theorem~\ref{thm:lower} in view. We also record the predicted mechanism for the eventual
formula, as a signpost rather than a claim: quaternionic $S$-spectrum theory
\cite{ColomboSabadiniStruppa2011,ACK2016} is natively conjugation-paired --- its resolvent kernel
$(T^2-2\re(s)T+|s|^2)^{-1}$ is a function of $(s,\bar s)$ by construction, and axially symmetric
zero sets of exactly our type are exactly the $S$-spectra --- so a trace formula for
$S$-resolvents would be a conjugation-paired explicit formula; to our knowledge that machinery has
never been aimed at number theory.

%=====================================================================
\section{Numerical certifications and reproducibility}\label{sec:numerics}

Every numerical claim in the paper is reproducible from the ancillary files
(\texttt{anc/}: \texttt{xi\_lib.py}, \texttt{explicit\_formula.py}, \texttt{si\_lib.py},
drivers \texttt{run\_validation.py}, \texttt{run\_experiments.py}, \texttt{run\_bridge.py},
\texttt{run\_si\_check.py}, the generator \texttt{make\_zeros.py}, a precomputed
\texttt{zeros.txt}, and a README mapping each script to the claims it certifies; Python with
numpy, scipy, mpmath, matplotlib; total runtime about $15$ minutes on commodity hardware).
Figures~\ref{fig:lab} and \ref{fig:bridge} are generated by the included scripts.

\emph{Validation} (\texttt{run\_validation.py}). The extension \eqref{eq:cosine} ($6$ theta
terms, trapezoid rule, $du=5\times10^{-5}$ on $[0,3.2]$) matches $\xi(\tfrac12+ix)$ computed
independently via mpmath to relative error $\le8\times10^{-10}$ at $x\in\{0,5,20,30\}$
(at $x\le20$ the agreement is $\sim10^{-14}$ or better), and vanishes at the first zero ordinate
to $7\times10^{-18}$, the quadrature floor. $\Xi_{DH}$ is real and even to relative $10^{-16}$
at $30$-digit precision (sampled); Newton iteration on $f_{DH}$ from crude seeds recovers the
off-line zeros \eqref{eq:DHspheres} with residuals $\le8\times10^{-30}$; the DH on-line
ordinates in $[74,97]$ are recovered by sign-change bisection.

\emph{Positivity experiments} (\texttt{run\_experiments.py}). Laguerre: relative margin of
$L_1[\Xi]$ has minimum $0.173$ on $[0,52]$ (beyond $x\approx55$ the float64 quadrature floor is
reached); for DH (rescaled $W=\eE^{\pi x/4}\Xi_{DH}$, sign-preserving since
$L_1[\eE^{cx}g]=\eE^{2cx}L_1[g]$), $L_1<0$ exactly on $[85.417,85.982]$ (predicted
$[85.391,86.008]$) and $[114.025,114.302]$ (predicted $[114.013,114.314]$), the shaving
accounted for by neighboring on-line zeros, while the sphere-free control $[78,82]$ has
$\min L_1\approx7.2\times10^2>0$. Sector energies (Gaussian weight, $20$ real ordinates plus an
implanted sphere at $a=28$ with $y_0\in[10^{-3},10^{-0.5}]$): log--log slopes
$2.000/4.000/6.000$ for $\ell=1,2,3$; $\EE_1$ of the real ordinates alone is exactly $0$;
per-sphere $\EE_1$ for the three DH spheres: $4.360\times10^{-3}$, $1.062\times10^{-3}$,
$2.593\times10^{-4}$ (successive ratios $4.11$, $4.10$ vs.\ predicted $(y_0/y_0')^2=4.18$,
$4.11$).

\emph{Bridge} (\texttt{run\_bridge.py}). Tables~\ref{tab:EF} and \ref{tab:rate}; the Gegenbauer
identity \eqref{eq:gegenbauer2} below holds to $\le3\times10^{-14}$ at $z\in\{0.5,1,2,4\}$ with
$80$ terms, where
\begin{equation}\label{eq:gegenbauer2}
\frac{\sinh(2z)}{2z}=\sum_{\ell\ge0}(2\ell+1)\,i_\ell(z)^2
\end{equation}
(the modified-Bessel decomposition of the sphere anomaly kernel over the same sectors as
\S\ref{sec:positivity}; specialize \cite[10.60.2, 10.47.12]{DLMF}, classically
\cite[\S11.4]{Watson1944}). The implant fit in Table~\ref{tab:rate} is biased because
$\eE^{2\lambda y_0}=\eE^5$ at the window edge barely clears the oscillatory background; the DH
sphere, with $\eE^{2\lambda y_0}=\eE^{15.4}$, gives the rate to four digits.

\begin{table}[ht]
\centering
\caption{Certification of \eqref{eq:EF} on the family \eqref{eq:Z}, $\sigma=2$, $250$ zeros:
zero side vs.\ pole${}+{}$archimedean${}+{}$prime side.}\label{tab:EF}
\begin{tabular}{lll}
\toprule
$(c,\lambda)$ & zero side & $|\Delta|$ \\
\midrule
$(20,0)$ & $1.86901176$ & $1.1\times10^{-10}$\\
$(20,1)$ & $0.96584078+1.42563531i$ & $1.2\times10^{-9}$\\
$(20,2)$ & $-0.85353747+1.53365468i$ & $6.4\times10^{-9}$\\
$(20,3)$ & $-1.80938739+0.21232192i$ & $1.8\times10^{-8}$\\
$(35,2)$ & $-0.38300486+0.14677138i$ & $4.8\times10^{-9}$\\
\bottomrule
\end{tabular}
\end{table}

\begin{table}[ht]
\centering
\caption{Rate experiment: $\sigma=2$, $\lambda\in[0.5,25]$, rate fitted to $\log|Z|$ on
$\lambda>15$.}\label{tab:rate}
\begin{tabular}{lccc}
\toprule
zero set & $|Z|^2$ at $\lambda=25$ & fitted rate & true $\sup y_\rho$ \\
\midrule
$\zeta$ zeros ($250$, real), $c=20$ & $3.2$ & $-0.0003$ & $0$\\
same $+$ implanted sphere $y_0=0.1$ at $a=c$ & $6.9\times10^2$ & $0.105$ & $0.1$\\
DH on-line ordinates only, $c=85.7$ & $5.8$ & $-0.009$ & $0$\\
DH with its first sphere, $c=85.7$ & $2.1\times10^7$ & $\mathbf{0.3085}$ & $\mathbf{0.30852}$\\
\bottomrule
\end{tabular}
\end{table}

\emph{Second involution} (\texttt{si\_lib.py}, \texttt{run\_si\_check.py}; $800$ zeros at $25$
digits, von Mangoldt data $n\le2\times10^4$, $\sigma=1$). The zero and prime sides of
$Z(c,\log2)$ agree to $2.7\times10^{-14}$ at $c\in\{30,100,250,500\}$;
$\osc_{[c-15,c+15]}F\ge2.01$ over $32$ windows spanning $[0,1160]$ and
$\osc=2.89,\,3.25,\,2.49$ at $c=10^4,10^5,10^6$ (prime side), against the guaranteed effective
bound $B_2/8=0.0676$ of Lemma~\ref{lem:osc}; and the encoding of Theorem~\ref{thm:encoding} is
demonstrated end to end on a synthetic multiset (two on-line ordinates near $100$ replaced by the
quadruple $\pm100.0745\pm0.25i$): the defect, $\sup_cD=0.125$, is recovered at $13$ test heights
with worst error $1.8\times10^{-15}$, with correct zeros in the defect-free regions, and with no
arithmetic input about the synthetic multiset used or even defined --- the numerical face of the
diagnosis of \S\ref{sec:resolution}. (The demonstration uses the comfortable windows
$[c-15,c+15]$, where the empirical oscillation constant is $r=2$, and a variant encoder; the
$|c_j-c|\le\sigma$ form of Theorem~\ref{thm:encoding} holds with a smaller $r$.)

\begin{figure}[p]
\centering
\includegraphics[width=\textwidth]{lab_results.pdf}
\caption{Left: $L_1[\Xi]>0$ on $[0,60]$ (semilog). Center: $L_1[\Xi_{DH}]$ dips negative
exactly inside the predicted windows $|x-a|<y_0$ at the first two DH spheres (shaded).
Right: sector energies $\EE_\ell\propto y_0^{2\ell}$, measured slopes $2.00/4.00/6.00$.}
\label{fig:lab}
\end{figure}

\begin{figure}[p]
\centering
\includegraphics[width=0.8\textwidth]{bridge_results.pdf}
\caption{The rate criterion (Theorem~\ref{thm:F}) on real data: $|Z(c,\lambda)|^2$ versus
$\lambda$ for $250$ real Riemann ordinates ($c=20$), the same plus an implanted sphere of radius
$0.1$, the DH on-line ordinates near $c=85.7$, and DH with its first sphere ($y_0=0.30852$;
fitted rate $0.3085$).}
\label{fig:bridge}
\end{figure}

\subsection*{Acknowledgements}
The synthesis of this program, the twistor-lift and explicit-formula computations, the numerical
experiments, and substantial portions of the exposition were carried out with Claude (Anthropic)
as a research assistant in August 2026. In particular, the Second Involution Conjecture was
formulated in that collaboration; when the assistant was subsequently asked to attempt a proof,
the outcome was the resolution reported in \S\S\ref{sec:resolution}--\ref{sec:repair} ---
Theorem~\ref{thm:encoding}, Lemma~\ref{lem:osc}, Theorem~\ref{thm:lower}, and the repaired
conjectures --- whose mathematics was drafted by the assistant and then checked by two further
independent AI reviewers (an adversarial proof audit and a blind re-derivation of the formulas and
constants), with the errors they found --- all in the first draft of the lower-bound section ---
repaired. All mathematical claims were reviewed by the author, who takes sole responsibility for
them, while not claiming the expertise to referee \S\S\ref{sec:resolution}--\ref{sec:repair};
those sections have not yet had expert human review, which is warmly invited. Code and archival materials are available at \cite{ZenodoArchive}.

%=====================================================================
\begin{thebibliography}{99}
\itemsep2pt

\bibitem{ACK2016}
D.~Alpay, F.~Colombo, and D.~P. Kimsey,
\emph{The spectral theorem for quaternionic unbounded normal operators based on the $S$-spectrum},
J. Math. Phys. \textbf{57} (2016), 023503.

\bibitem{Altavilla2018}
A.~Altavilla,
\emph{Twistor interpretation of slice regular functions},
J. Geom. Phys. \textbf{123} (2018), 184--208.

\bibitem{AltavillaSarfatti2019}
A.~Altavilla and G.~Sarfatti,
\emph{Slice-polynomial functions and twistor geometry of ruled surfaces in $\mathbb{CP}^3$},
Math. Z. \textbf{291} (2019), 1059--1092.

\bibitem{Baker2015}
E.~B. Baker III,
\emph{Quaternionic analysis with applications to the Riemann zeta function},
manuscript (2015); archived in \cite{ZenodoArchive}.

\bibitem{Baker2016}
E.~B. Baker III,
\emph{The quaternion-regular exponential and the three dimensional Clifford--Fourier kernel},
manuscript (2016); archived in \cite{ZenodoArchive}.

\bibitem{Baker2017}
E.~B. Baker III,
\emph{Twistor holomorphicity},
manuscript (2017); archived in \cite{ZenodoArchive}.

\bibitem{Baker2022arXiv}
E.~B. Baker III,
\emph{Causal and self-dual morphisms in four complex dimensions},
arXiv:2203.07952 (2022).

\bibitem{Baker2023arXiv}
E.~B. Baker III,
\emph{Generalized conformal maps as classical symmetries of Yang--Mills fields},
arXiv:2301.00029 (2023).

\bibitem{ZenodoArchive}
E.~B. Baker III,
\emph{The zero-sphere program: archival manuscripts, working memos, and code (2015--2026)},
Zenodo, DOI: [to appear].

\bibitem{BCR}
J.~Bochnak, M.~Coste, and M.-F. Roy,
\emph{Real Algebraic Geometry},
Ergeb. Math. Grenzgeb. (3) \textbf{36}, Springer, Berlin, 1998.

\bibitem{Bombieri2000}
E.~Bombieri,
\emph{Remarks on Weil's quadratic functional in the theory of prime numbers, I},
Atti Accad. Naz. Lincei Rend. Lincei Mat. Appl. \textbf{11} (2000), 183--233.

\bibitem{CGSS2009}
F.~Colombo, G.~Gentili, I.~Sabadini, and D.~C. Struppa,
\emph{Extension results for slice regular functions of a quaternionic variable},
Adv. Math. \textbf{222} (2009), 1793--1808.

\bibitem{ColomboSabadiniStruppa2011}
F.~Colombo, I.~Sabadini, and D.~C. Struppa,
\emph{Noncommutative Functional Calculus},
Progress in Mathematics \textbf{289}, Birkh\"auser, Basel, 2011.

\bibitem{Connes2016}
A.~Connes,
\emph{An essay on the Riemann Hypothesis},
in \emph{Open Problems in Mathematics} (J.~F. Nash Jr., M.~Th. Rassias, eds.), Springer, 2016,
225--257.

\bibitem{CsordasVarga1990}
G.~Csordas and R.~S. Varga,
\emph{Necessary and sufficient conditions and the Riemann hypothesis},
Adv. in Appl. Math. \textbf{11} (1990), 328--357.

\bibitem{CsordasVishnyakova2013}
G.~Csordas and A.~Vishnyakova,
\emph{The generalized Laguerre inequalities and functions in the Laguerre--P\'olya class},
Cent. Eur. J. Math. \textbf{11} (2013), 1643--1650.

\bibitem{Davenport2000}
H.~Davenport,
\emph{Multiplicative Number Theory}, 3rd ed., revised by H.~L. Montgomery,
Graduate Texts in Mathematics \textbf{74}, Springer, New York, 2000.

\bibitem{DavenportHeilbronn1936a}
H.~Davenport and H.~Heilbronn,
\emph{On the zeros of certain Dirichlet series (I)},
J. London Math. Soc. \textbf{11} (1936), 181--185.

\bibitem{DavenportHeilbronn1936b}
H.~Davenport and H.~Heilbronn,
\emph{On the zeros of certain Dirichlet series (II)},
J. London Math. Soc. \textbf{11} (1936), 307--312.

\bibitem{DLMF}
\emph{NIST Digital Library of Mathematical Functions}, \url{https://dlmf.nist.gov/},
Release 1.2.x; F.~W.~J. Olver et al., eds.

\bibitem{GentiliStoppato2008}
G.~Gentili and C.~Stoppato,
\emph{Zeros of regular functions and polynomials of a quaternionic variable},
Michigan Math. J. \textbf{56} (2008), 655--667.

\bibitem{GSSbook}
G.~Gentili, C.~Stoppato, and D.~C. Struppa,
\emph{Regular Functions of a Quaternionic Variable}, 2nd ed.,
Springer Monographs in Mathematics, Springer, Cham, 2022.

\bibitem{GentiliStruppa2007}
G.~Gentili and D.~C. Struppa,
\emph{A new theory of regular functions of a quaternionic variable},
Adv. Math. \textbf{216} (2007), 279--301.

\bibitem{GSS2014}
G.~Gentili, S.~Salamon, and C.~Stoppato,
\emph{Twistor transforms of quaternionic functions and orthogonal complex structures},
J. Eur. Math. Soc. (JEMS) \textbf{16} (2014), 2323--2353.

\bibitem{GhiloniPerotti2011}
R.~Ghiloni and A.~Perotti,
\emph{Slice regular functions on real alternative algebras},
Adv. Math. \textbf{226} (2011), 1662--1691.

\bibitem{GoldstonMontgomery1987}
D.~A. Goldston and H.~L. Montgomery,
\emph{Pair correlation of zeros and primes in short intervals},
in \emph{Analytic Number Theory and Diophantine Problems}, Birkh\"auser, 1987, 183--203.

\bibitem{GriffithsHarris1978}
P.~Griffiths and J.~Harris,
\emph{Principles of Algebraic Geometry},
Wiley, New York, 1978.

\bibitem{GriffithsHarris1979}
P.~Griffiths and J.~Harris,
\emph{Algebraic geometry and local differential geometry},
Ann. Sci. \'Ecole Norm. Sup. (4) \textbf{12} (1979), 355--432.

\bibitem{Guinand1948}
A.~P. Guinand,
\emph{A summation formula in the theory of prime numbers},
Proc. London Math. Soc. (2) \textbf{50} (1948), 107--119.

\bibitem{IwaniecKowalski2004}
H.~Iwaniec and E.~Kowalski,
\emph{Analytic Number Theory},
AMS Colloquium Publications \textbf{53}, Amer. Math. Soc., Providence, RI, 2004.

\bibitem{Montgomery1973}
H.~L. Montgomery,
\emph{The pair correlation of zeros of the zeta function},
Proc. Sympos. Pure Math. \textbf{24}, Amer. Math. Soc., 1973, 181--193.

\bibitem{Polya1926}
G.~P\'olya,
\emph{Bemerkung \"uber die Integraldarstellung der Riemannschen $\xi$-Funktion},
Acta Math. \textbf{48} (1926), 305--317.

\bibitem{RodgersTao2020}
B.~Rodgers and T.~Tao,
\emph{The de Bruijn--Newman constant is non-negative},
Forum Math. Pi \textbf{8} (2020), e6.

\bibitem{Spira1994}
R.~Spira,
\emph{Some zeros of the Titchmarsh counterexample},
Math. Comp. \textbf{63} (1994), 747--748.

\bibitem{Titchmarsh1986}
E.~C. Titchmarsh,
\emph{The Theory of the Riemann Zeta-Function}, 2nd ed., revised by D.~R. Heath-Brown,
Oxford University Press, 1986.

\bibitem{Watson1944}
G.~N. Watson,
\emph{A Treatise on the Theory of Bessel Functions}, 2nd ed.,
Cambridge University Press, 1944.

\bibitem{Weil1952}
A.~Weil,
\emph{Sur les ``formules explicites'' de la th\'eorie des nombres premiers},
Comm. S\'em. Math. Univ. Lund (1952), Tome Suppl\'ementaire, 252--265.

\bibitem{Zak1993}
F.~L. Zak,
\emph{Tangents and Secants of Algebraic Varieties},
Translations of Mathematical Monographs \textbf{127},
Amer. Math. Soc., Providence, RI, 1993.

\end{thebibliography}

\end{document}
